\chapter{The honest ledger}
\label{ch:ledger}

\section*{What this chapter is for}

The most counter-intuitive move available in a portfolio is to write down, in
the repository, a dated list of everything it does not do. It costs an hour, it
reads as seniority rather than as weakness, and it has a second effect that
matters more: it stops you saying something in an interview that your own work
cannot support.

\section{The register}

A deviations register is a file at a known path listing, for each thing the
artefact does not do or does not do properly:

\begin{itemize}
  \item what it is,
  \item why it is that way,
  \item the date the decision was taken,
  \item and \textbf{what would close it}.
\end{itemize}

The fourth field does the work. Without it the file is a list of excuses; with
it, every row is a decision with an exit condition, which is what an engineering
decision looks like when somebody is being serious.

\begin{kitnote}
\textbf{An acknowledged deviation is a decision. An unacknowledged one is
drift.}

That sentence is the whole discipline. It also gives you the rule for
maintaining the file: when a row is fixed, \emph{delete} it. Do not mark it
resolved. A register that only grows is a register nobody reads, and a row that
stays open indefinitely teaches the reader that the principle it violates was
optional.
\end{kitnote}

\section{Why it reads as seniority}

Three reasons, and the third is the one that decides interviews.

\textbf{It is the difference between a demo and a system.} Anybody can build
something that works on the path they walked. Knowing which paths you did not
walk requires having thought about the whole surface.

\textbf{It converts every awkward question into a prepared one.} An interviewer
finding a hole they think you missed is a bad moment. An interviewer finding a
hole listed on line 14 with a date and a closing condition is a different
conversation entirely, and the thing being demonstrated is no longer the hole.

\textbf{It is the only credible way to claim rigour.} A repository that claims
to be careful and lists nothing undone is making a claim no experienced reader
believes, because nothing is complete. \judgement{The register is what makes
the rest of the README credible; without it, the careful language reads as
marketing.}

\section{What belongs in it}

\begin{kittable}{@{}lX@{}}
\textbf{Class} & \textbf{Example} \\
\midrule
Never run & an integration written against a published specification and never executed against a live one \\
Run once & a path exercised manually and not covered by a check \\
Untested by construction & behaviour that cannot be tested without a credential you do not have \\
Deliberately out of scope & something a real version needs and this one does not attempt \\
A standard you did not meet & where you knowingly departed from your own convention and why \\
Blocked on something outside the work & a platform feature, an account, an approval \\
\end{kittable}

The first and third are the valuable ones, because they are the ones a reader
cannot discover from the code. A module that has never been executed against
the real thing compiles, passes its unit tests, and is an unverified claim.

\section{The pattern this generalises}

Once you have written one of these, you will notice it applies to more than a
repository.

\begin{kitmethod}
\textbf{The artefact} gets a deviations register.

\textbf{Your application documents} get a claim inventory
(Chapter~\ref{ch:claims}) \dash{} the same idea, applied to sentences instead of
code.

\textbf{You, before a call,} get a banned-claims list: the things your evidence
does not support, written down so you do not say them under pressure.

\textbf{This book} gets Appendix~\ref{app:refused}, which names the figures it
refuses to repeat and why.
\end{kitmethod}

\judgement{These are one habit, not four. The habit is that a claim you have
not checked is tracked as an open item rather than quietly carried, and the
reason it transfers so well into hiring is that a hiring process is almost
entirely an exchange of claims.}

\section{The failure mode}

A register can be maintained by not looking. A short list is the goal; an empty
one achieved by never auditing is a different thing wearing the same clothes.

So the file carries the date it was last checked against the repository, and the
check is a re-read, not a memory. \judgement{If you cannot remember the last
time you opened it, the dates in it are decorations.}

\begin{drill}
Take the artefact you would send to an employer today. Set a timer for forty
minutes and write its register: every integration you never ran live, every
path you tested by hand, every standard of your own you did not meet. Put a
date on it.

Then read the list and find the row you would least like to be asked about.
That row is the first thing to fix, and until it is fixed it is the first thing
to say out loud.
\end{drill}

\section*{What this chapter does not claim}

It does not claim every interviewer will read the register or reward it. Some
will not open the repository at all. The claim is that the register is cheap,
that it changes what you are able to say honestly, and that the second effect
is worth the hour even in the loops where nobody reads the file.

The specific position that honesty about gaps reads as seniority is a
judgement, and it is the author's. It is stated as one.
